%=============================================================
% Paper 12: Hopenhayn, Neira & Singhania (2022), Econometrica
%=============================================================
\chapter{Hopenhayn, Neira \& Singhania (2022)}

\begin{paperbox}{From Population Growth to Firm Demographics: Implications
for Concentration, Entrepreneurship and the Labor Share}
\textbf{Authors:} Hugo Hopenhayn (UCLA), Julian Neira (Exeter),
Rish Singhania (Exeter)\\[3pt]
\textbf{Journal:} \textit{Econometrica} 90(4): 1879--1914\\[3pt]
\textbf{Year:} 2022 \quad\textbf{Field:} Macroeconomics; Firm dynamics;
Labour economics\\[3pt]
\textbf{Classification:} Structural (GE model + quantitative calibration
+ historical decomposition)
\end{paperbox}

%------------------------------------------------------------
\section{Research Question}

Can changes in population/labour force growth---primarily driven by the baby
boom and subsequent demographic slowdown---provide a unified quantitative
explanation for three major secular trends in the U.S.\ economy: (i)~rising
concentration of employment in large firms; (ii)~declining firm entry rates;
and (iii)~declining aggregate labour share?

%------------------------------------------------------------
\section{Motivation}

Three well-documented but apparently disconnected secular trends in the U.S.\
economy since the late 1970s have attracted separate literatures:
\begin{itemize}
  \item \textbf{Rising concentration:} employment share in firms with
    250+ employees rose from 51.6\% to 57.3\% (1978--2014).
  \item \textbf{Declining dynamism:} U.S.\ firm entry rate fell by
    approximately 6 percentage points.
  \item \textbf{Declining labour share:} aggregate share of GDP going
    to labour has fallen substantially.
\end{itemize}
The paper proposes a single unified mechanism: the baby boom increased firm
entry rates in the 1970s, creating a cohort of ``baby boom firms''; as these
firms aged, the firm-age distribution shifted toward older, larger, less dynamic
firms---mechanically generating all three trends simultaneously through a simple
accounting identity.

%------------------------------------------------------------
\section{Main Contribution}

\begin{enumerate}
  \item \textbf{Empirical documentation:} Using BDS data (1978--2014),
    shows that once firm-age controls are included, the time trends in
    concentration, exit rates, and average size \emph{reverse sign}---the
    aggregate trends are entirely driven by changes in the firm-age
    distribution.

  \item \textbf{Invariance theorem (Corollary~1):} In a general class of
    GE models (perfect competition, CES monopolistic competition, variable
    markup oligopoly), equilibrium allocations \emph{conditional on firm
    age are invariant to changes in population growth}. All aggregate
    changes must therefore come from changes in the firm-age distribution.

  \item \textbf{Dynamic entry equation (Corollary~2):} Characterises
    firm entry as a history-dependent distributed-lag equation, enabling
    quantitative simulation from historical labour force data.

  \item \textbf{Quantitative analysis:} Feeds historical U.S.\ labour
    force data (1940--2014) through the model; replicates the observed
    entry rate decline, concentration, average size, exit rates, and the
    aggregate labour share (including the baby-boom spike and subsequent
    decline).
\end{enumerate}

%------------------------------------------------------------
\section{Relationship to the Literature}

Synthesises three previously disconnected literatures: (i)~business
dynamism (Haltiwanger, Jarmin, Miranda 2011; Decker et al.\ 2014, 2020);
(ii)~rising concentration (Autor et al.\ 2020; Grullon et al.\ 2019);
(iii)~labour share decline (Karabarbounis--Neiman 2014; Kehrig--Vincent
2021). The closest predecessor is Karahan, Pugsley, \c{S}ahin (2018), who
link population growth to entry rates in a Hopenhayn (1992)-style model.
The present paper extends this by providing conditions applicable to imperfect
competition, characterising transitional dynamics explicitly, and showing
that baby-boom transitional effects explain 55\% of the entry rate decline
--- more than the long-run effect (45\%).

%------------------------------------------------------------
\section{Theoretical Framework and Economic Mechanism}

\subsection*{Accounting Identity}

The fundamental identity of the paper:
\begin{equation}
  \lambda_t = \frac{\dot{N}_t}{N_{t-1}}\frac{e_{t-1}}{e_t} - 1 + \xi_t,
  \label{eq:hns_identity}
\end{equation}
where $\lambda_t$ is the firm entry rate, $\dot{N}/N$ is labour force growth,
$e_t = N_t/M_t$ is average employment per firm, and $\xi_t$ is the aggregate
exit rate. Entry rate equals labour force growth minus growth in average size
plus exit rate.

\subsection*{General Framework}

Labour endowment $N_t$ inelastically supplied. Firm profits
$\pi(s,z)$ and employment $n(s,z)$ are increasing in idiosyncratic state~$s$
and aggregate state~$z$, under a Markov process $F(s'|s)$. The firm's
Bellman equation:
\begin{equation}
  v(s,z_t) = \max\!\left\{0,\;\pi(s,z_t)+\beta\,\mathbb{E}[v(s',z_{t+1})|s]\right\},
  \label{eq:hns_bellman}
\end{equation}
with exit threshold $s^*_t = \inf\{s\,|\,\pi(s,z_t)+\beta\mathbb{E}[v(s',z_{t+1})|s]>0\}$.
Free entry pins down $z^*$ via $v^e(z^*) = \int v(s,z^*)\,dG(s) - c_e = 0$.

\subsection*{Invariance of Age Profiles (Corollary~1)}

Because the aggregate state $z^*$ is determined solely by the zero-profit
condition for entrants---independent of $N_t$---adjustments in aggregate
quantities occur entirely through changes in the mass of entrants $m_t$,
leaving age profiles unchanged:
\begin{align*}
  &\text{Exit rates by age } \tilde{\xi}_a, \quad
  \text{Average firm size by age } \tilde{e}_a, \quad
  \text{Size distribution by age}\\
  &\hspace{2cm}\text{are all invariant to changes in } N_t.
\end{align*}
This result holds for all models satisfying the general framework
(three examples given: perfect competition, CES monopolistic competition,
variable markup oligopoly).

\subsection*{Dynamic Entry Equation (Corollary~2)}

The invariance of age profiles implies that market clearing
$N_t = m_t(S_0\tilde{e}_0+c_e) + \sum_{a=1}^{t}m_{t-a}S_a\tilde{e}_a$
uniquely determines:
\begin{equation}
  m_t = \frac{N_t - \sum_{a=1}^{t}m_{t-a}\,S_a\,\tilde{e}_a}{S_0\,\tilde{e}_0 + c_e},
  \label{eq:hns_entry}
\end{equation}
where $S_a$ is the probability of surviving to age~$a$ and $\tilde{e}_a$
is average firm size at age~$a$. Entry is \emph{history-dependent}: current
entry depends on the entire distributed lag of past entry masses
$\{m_{t-a}\}$.

\subsection*{Long-Run Multiplier (Proposition~2)}

When $\tilde{\xi}_a$ is decreasing in $a$ (older firms have lower exit
rates, as in the data), the long-run aggregate exit rate $\xi^{SS}$ is
\emph{increasing} in population growth~$g$. Therefore, the effect of
population growth on the entry rate is amplified:
\begin{equation}
  \lambda^{SS} = g + \xi^{SS} \quad\Rightarrow\quad
  \frac{d\lambda^{SS}}{dg} = 1 + \frac{d\xi^{SS}}{dg} \approx 1.5
  \label{eq:hns_multiplier}
\end{equation}
for the U.S.\ economy. A 2~pp decline in labour force growth generates
$\approx 3$~pp decline in the long-run entry rate (long-run effect alone).

\subsection*{Baby Boom Transitional Dynamics}

The baby boom raised firm entry rates in the 1960s--70s, creating a glut
of firms. Three reinforcing feedback loops amplify the direct effect of
declining population growth:
\begin{enumerate}
  \item \textbf{Exit rate feedback:} Aging distribution lowers aggregate
    exit rate $\xi_t$ (older firms exit less) $\Rightarrow$ entry falls
    further.
  \item \textbf{Average size feedback:} Aging distribution raises average
    firm size $e_t$ (older firms are larger) $\Rightarrow$ entry falls
    further.
  \item \textbf{Baby boom glut:} The large cohort of baby-boom firms has
    not yet died off; as of 2014, the firm-age distribution is older than
    the new steady state would imply.
\end{enumerate}

%------------------------------------------------------------
\section{Data}

\begin{itemize}
  \item \textbf{Business Dynamics Statistics (BDS):} U.S.\ Census Bureau,
    1978--2014; near-universal coverage of private-sector firms with paid
    employees; entry/exit rates, average firm size, concentration, and
    job flows by firm age and sector.
  \item \textbf{Historical labour force data:} 1940--2014 (BLS civilian
    labour force; Survey of Current Business entry rates 1940--1962).
  \item \textbf{Structural calibration:} GMM targeting 1978--1983 moments:
    entry rate (13.14\%), average entrant size (5.96), concentration of
    entrants, 5-year firm growth (73.86\%), 5-year exit rate (55.01\%),
    large-firm employment shares.
  \item \textbf{Labour share:} Karabarbounis--Neiman (2014); Koh,
    Santaeul{\`a}lia-Llopis, Zheng (2020).
\end{itemize}

%------------------------------------------------------------
\section{Empirical Strategy}

Three complementary approaches:
\begin{enumerate}
  \item \textbf{Regression evidence (Table~I):} Regress exit rate, average
    firm size, and concentration on year trend, with and without firm-age
    controls. The reversal of sign when age controls are included is the
    core identification.
  \item \textbf{Shift-share accounting:} Compute
    \[
      \text{Change if Aging Only}
      = \frac{(\text{Unconditional trend} - \text{Age-conditional trend})\times 36}
             {\text{Level}_{1978}}
    \]
    to isolate the contribution of the changing age distribution.
  \item \textbf{Dynamic entry equation simulation:} Estimate the firm
    employment stochastic process by GMM; feed historical labour force
    data from 1940 into equation~(\ref{eq:hns_entry}) and iterate forward
    to generate aggregate time series.
\end{enumerate}

%------------------------------------------------------------
\section{Identification Strategy}

The key identifying assumption is the \textbf{invariance of age profiles}:
exit rates, average size, and size distributions by firm age do not change
with population growth. This is directly validated by the regression evidence:
\begin{itemize}
  \item After controlling for firm age, the exit rate trend becomes slightly
    \emph{positive} (consistent with no within-age-group change).
  \item After controlling for firm age, average size and concentration
    trends become \emph{negative} (consistent with mild within-age
    declines).
\end{itemize}
The aggregate positive trends in size and concentration are entirely accounted
for by the shifting firm-age distribution.

\textbf{Important caveat:} The invariance result requires a perfectly elastic
supply of potential entrants proportional to population. If the pool of
potential entrepreneurs is fixed, the mechanism is weaker.

%------------------------------------------------------------
\section{Main Results}

\subsection*{Regression Evidence (Table~I)}

\begin{center}
\small
\begin{tabular}{lccc}
\toprule
\textbf{Variable} & \textbf{Aggregate trend} &
  \textbf{Controlling for age} & \textbf{Sign reversal?} \\
\midrule
Exit rate & $-0.043$/yr *** & $+0.011$/yr & Yes \\
Average firm size & $+0.096$/yr *** & $-0.143$/yr *** & Yes \\
Concentration & $+0.159$/yr *** & $-0.069$/yr *** & Yes \\
\bottomrule
\end{tabular}
\end{center}

\subsection*{Decomposition of 6~pp Entry Rate Decline (Table~V)}

\begin{center}
\begin{tabular}{lcc}
\toprule
\textbf{Component} & \textbf{Magnitude} & \textbf{Share} \\
\midrule
Long-run effect (steady state) & $-2.76$~pp & 45\% \\
1978 transition effect (baby boom entry) & $-1.47$~pp & 24\% \\
2014 transition effect (baby boom aging) & $-1.84$~pp & 30\% \\
Residual & $-0.01$~pp & 1\% \\
\midrule
\textbf{Total} & $-6.08$~pp & 100\% \\
\bottomrule
\end{tabular}
\end{center}

Of the long-run effect: direct labour force growth decline accounts for
1.88~pp; long-run multiplier effect on exit rates adds 0.88~pp.

\subsection*{Model Fit for Aggregate Variables (Figure~8)}

Aging alone (model) generates:
\begin{itemize}
  \item 16\% decline in exit rate (data ``aging only'': 18\%)
  \item 41\% increase in average firm size (data: 43\%)
  \item 15\% increase in concentration (data: 16\%)
\end{itemize}

\subsection*{Labour Share and Job Reallocation}

Firm aging generates the hump-shaped post-WWII pattern in the aggregate
labour share: rising until $\approx$1980 (baby-boom entrants are small,
high-labour-share) and falling after (firms age and grow). Firm aging
explains:
\begin{itemize}
  \item \textbf{48\%} of the decline in the job creation rate
  \item \textbf{38\%} of the decline in the job destruction rate
  \item \textbf{32\%} of the decline in the job reallocation rate
\end{itemize}

%------------------------------------------------------------
\section{Economic Magnitude and Interpretation}

The 6~pp decline in the U.S.\ firm entry rate represents a massive
structural change. The paper shows that baby-boom transitional dynamics
(55\% of total) are quantitatively larger than the long-run steady-state
effect (45\%), meaning the timing of the demographic cycle matters as much
as its level. With labour force projections through 2060, entry rates are
projected to partially recover by $\approx$2030 as the baby boom cohort
exits and millennial births enter the labour force.

%------------------------------------------------------------
\section{Mechanisms}

Three reinforcing feedback loops amplify the direct demographic effect:
\begin{enumerate}
  \item \textbf{Exit rate feedback:} Declining entry $\to$ aging
    distribution $\to$ lower aggregate exit rate (older firms exit less)
    $\to$ even lower entry.
  \item \textbf{Average size feedback:} Aging distribution $\to$ higher
    average firm size (older firms are larger) $\to$ lower entry demand
    per unit of labour.
  \item \textbf{Baby boom glut:} The large cohort of firms created in the
    1970s continues to suppress entry in 2014 because the 2014 firm-age
    distribution is older than the new steady state would predict.
\end{enumerate}

%------------------------------------------------------------
\section{Robustness Checks}

\begin{itemize}
  \item Results hold in all 7 major BDS sectors (Table~II); aging is
    quantitatively important across every sector.
  \item Sector controls do not remove aggregate trends; they deepen them
    for size and concentration.
  \item Invariance assumption validated by regression evidence (age controls
    flip sign of all three trends).
  \item Shift-share exercise using data levels and model age weights yields
    similar results.
  \item Alternative explanations evaluated within the model: higher entry
    costs, economies of scale, lower diffusion rates---each can explain some
    facts but is inconsistent with others.
  \item WW2 fluctuations: model generates entry rate fluctuations matching
    Survey of Current Business data for 1940--1962.
\end{itemize}

%------------------------------------------------------------
\section{Limitations}

\begin{enumerate}
  \item \textbf{Perfectly elastic entrant supply:} Requires the number of
    potential entrants proportional to population; weakened if the
    entrepreneurial pool is fixed.
  \item \textbf{Constant aggregate state:} Invariance result requires
    risk-neutral households or small open economy; CRRA with endogenous
    interest rates may partially offset the mechanism.
  \item \textbf{No welfare analysis:} The paper does not characterise
    whether the shift to older, larger firms is efficient or inefficient.
  \item \textbf{Labour share mechanism:} Requires overhead labour (fixed
    cost) to generate the size--labour-share relationship; maintained
    assumption rather than derived from micro-foundations.
\end{enumerate}

%------------------------------------------------------------
\section{Policy Implications}

\begin{itemize}
  \item \textbf{Declining entrepreneurship is a demographic phenomenon.}
    Policies targeting startup subsidies address the symptom; the root cause
    (declining birth rates) will not be reversed by business policy.
  \item \textbf{Concentration is not necessarily a competition problem.}
    Rising concentration reflects a shift in the age distribution of firms,
    not necessarily market power. This complicates antitrust interpretations.
  \item \textbf{Labour share decline is partly demographic.} Interventions
    targeting labour market institutions or trade may address only a subset
    of the decline.
  \item \textbf{Job reallocation decline.} 32--48\% of the decline in gross
    job reallocation is explained by firm aging; restoring dynamism requires
    addressing the demographic slowdown, not just regulatory reform.
\end{itemize}

%------------------------------------------------------------
\section{Overall Assessment}

Hopenhayn, Neira, and Singhania (2022) achieve something rare in
macroeconomics: a parsimonious unified explanation for three major trends
that have separately occupied large literatures. The invariance result,
applicable across a general class of GE models, is a rigorous theoretical
contribution. The transition dynamics analysis---distinguishing long-run
and baby-boom transitional effects---is both analytically and quantitatively
careful. The empirical evidence (sign reversal with age controls) is
compelling.

\textbf{Quality: Outstanding. A landmark paper in firm dynamics and
macroeconomics.}

%------------------------------------------------------------
\section{Relevance to My Research}

Highly relevant for RDC research on Canadian firm dynamics. Canada has
experienced similar demographic trends. The Business Register/LEAP data
could test the HNS mechanism for Canada: does firm aging explain declining
Canadian entry rates? The heterogeneous provincial demographics (Quebec
birth rate fluctuations, Alberta immigration-driven growth) could provide
within-country identifying variation.

%------------------------------------------------------------
\section{Potential Research Extensions}

\begin{itemize}
  \item Test the HNS mechanism for Canada using LEAP data; exploit
    provincial variation in demographic patterns (Quebec vs.\ Alberta
    vs.\ Ontario) as quasi-experimental variation in the key driving force.
  \item Extend the model to incorporate immigration separately from
    native birth rates; test whether immigration-driven labour force growth
    has different effects on firm demographics than birth-rate-driven growth.
\end{itemize}

%------------------------------------------------------------
\section{Research Gaps}

\begin{itemize}
  \item The welfare consequences of the baby-boom firm aging are
    uncharacterised: is the shift toward older, larger firms efficient
    (positive selection) or inefficient (protection of incumbents from
    younger challengers)?
  \item The model does not account for the global nature of the trends;
    similar patterns in European economies with different demographics
    suggest additional forces beyond demographics alone.
\end{itemize}

%------------------------------------------------------------
\section{Methodological Lessons}

\begin{enumerate}
  \item The accounting identity (\ref{eq:hns_identity}) organises all
    subsequent analysis without requiring a full structural model; it
    enables clear decomposition of what drives entry rate changes.
  \item The regression reversal (age controls flip all three trend signs)
    is a simple, convincing empirical strategy for establishing the proximate
    mechanism, regardless of the underlying model.
  \item The invariance result reduces a complex GE problem to a demographic
    accounting exercise, making the dynamic entry equation (\ref{eq:hns_entry})
    applicable across market structures.
  \item Separating transitional from long-run effects reveals that one-half of
    the observed trend is transitory (baby boom)---the economy is not in steady
    state and may partially recover, with major implications for policy urgency.
\end{enumerate}

%------------------------------------------------------------
\section*{Five-Sentence Summary}

Hopenhayn, Neira, and Singhania (2022) document that three major secular
trends in the U.S.\ economy---rising employment concentration, declining firm
entry rates ($-6$~pp from 1978--2014), and declining labour share---all
reverse sign once firm-age controls are included in regressions, establishing
that they are driven entirely by changes in the firm-age distribution rather
than within-age-group trends. They develop a general theoretical framework,
applicable to perfect competition, CES monopolistic competition, and variable
markup oligopoly, showing that equilibrium allocations conditional on firm age
are invariant to changes in population growth, and derive a dynamic entry
equation linking current entry to the entire distributed lag of past entry via
survival probabilities and average firm size by age. Feeding historical U.S.\
labour force data (1940--2014) into this equation, the model replicates the
6~pp decline in the entry rate---with 45\% attributable to the long-run
multiplier effect of declining labour force growth (elasticity $\approx 1.5$)
and 55\% to baby-boom transitional dynamics (aging of firms founded in the
1970s). Firm aging also generates the hump-shaped post-WWII pattern in the
aggregate labour share and accounts for 32--48\% of the decline in gross job
reallocation rates. The paper provides the first unified quantitative
explanation linking baby-boom demographics to declining entrepreneurship,
rising concentration, declining labour share, and falling job reallocation---
trends previously treated as separate phenomena requiring separate explanations.

\vspace{6pt}
\begin{quickref}
\begin{tabularx}{\linewidth}{@{}lX@{}}
\textbf{Key contribution} &
  First unified quantitative explanation linking baby-boom demographics to
  declining entrepreneurship, rising concentration, declining labour share,
  and falling job reallocation through firm aging; invariance theorem
  applicable across a general class of GE models. \\[4pt]
\textbf{Key weakness} &
  Requires perfectly elastic entrant supply proportional to population; no
  welfare analysis of the aging transition; labour share mechanism requires
  maintained assumption of overhead labour. \\[4pt]
\textbf{Most interesting gap} &
  The welfare consequences of baby-boom firm aging are uncharacterised: is
  the shift toward older, larger firms efficient (positive selection) or
  inefficient (protection of incumbents from younger challengers)? \\[4pt]
\textbf{Publishable extension} &
  Test the HNS mechanism for Canada using LEAP data, exploiting provincial
  variation in demographic patterns (Quebec vs.\ Alberta vs.\ Ontario) as
  quasi-experimental variation; test whether the long-run multiplier
  elasticity of $\approx 1.5$ also holds in Canadian provincial data.
\end{tabularx}
\end{quickref}
